%======================================================================
\chapter{Problem Formulation}
\label{ch:problem}
%======================================================================

%----------------------------------------------------------------------
\section{Domain and Geometry}

The tissue domain $\Omega \subset \mathbb{R}^2$ is a square of side
$L_x \times L_y = \SI{5}{\cm} \times \SI{5}{\cm}$, discretized on a
uniform $N_x \times N_y$ grid.
Three non-overlapping sub-regions partition $\Omega$:
\begin{align}
    \Omega_T &= \{\mathbf{r} : \|\mathbf{r} - \mathbf{r}_c\| \leq r_T\}
        & &\text{(tumor)} \\
    \Omega_M &= \{\mathbf{r} : r_T < \|\mathbf{r} - \mathbf{r}_c\|
        \leq r_T + \delta_M\}
        & &\text{(safety margin)} \\
    \Omega_H &= \Omega \setminus (\Omega_T \cup \Omega_M)
        & &\text{(healthy tissue)}
\end{align}
where $\mathbf{r}_c$ is the tumor center, $r_T = \SI{8}{\mm}$ is the tumor
radius, and $\delta_M = \SI{3}{\mm}$ is the margin width.

%----------------------------------------------------------------------
\section{State Equations}

\subsection{State Equation 1 — Pennes Bioheat PDE}

The temperature field $T(\mathbf{r}, t)$ in \si{\degreeCelsius} satisfies:
\begin{equation}
    \rho(\mathbf{r})\, c(\mathbf{r})\,\frac{\partial T}{\partial t}
    = \nabla \cdot \bigl(k(\mathbf{r})\,\nabla T\bigr)
    - \omega_b \rho_b c_b \bigl(T - T_b\bigr)
    + Q_{\text{met}}(\mathbf{r})
    + Q_{\text{source}}(\mathbf{r},\, u,\, t)
    \label{eq:bioheat}
\end{equation}
with initial condition $T(\mathbf{r}, 0) = T_{\text{init}}$ and boundary
condition on $\partial\Omega$ (Dirichlet, Neumann, or Robin; see
Section~\ref{sec:bc}).

The volumetric heat source is modeled by a Gaussian near-field \gls{sar}
approximation:
\begin{equation}
    Q_{\text{source}}(\mathbf{r}, P, t)
    = \underbrace{\mathrm{SAR}_{\text{peak}}\,
        \exp\!\left(-\frac{\|\mathbf{r} - \mathbf{r}_p\|^2}{2\sigma_{\mathrm{SAR}}^2}\right)
      }_{\mathrm{SAR}(\mathbf{r})}\cdot P(t)
    \label{eq:sar}
\end{equation}
where $\mathbf{r}_p$ is the probe position, $\sigma_{\mathrm{SAR}}$ is the
beam spread, and $P(t)$ is the control input (applied power in watts).

After finite-difference discretization with $N = N_x N_y$ voxels:
\begin{equation}
    \dot{\mathbf{x}}_T
    = M^{-1}\!\left[K_d\,\mathbf{x}_T
      - W_b\bigl(\mathbf{x}_T - T_b\,\mathbf{1}\bigr)
      + \mathbf{q}_{\mathrm{met}}
      + \mathbf{b}_P\,P(t)\right]
    \label{eq:bioheat-discrete}
\end{equation}
where $M = \mathrm{diag}(\rho_i c_i)$, $K_d$ is the sparse FD Laplacian
weighted by $k$, $W_b = \mathrm{diag}(\omega_b \rho_b c_b)$, and
$\mathbf{b}_P = \mathrm{SAR}(\mathbf{r})$ vectorized.

\subsection{State Equation 2 — Arrhenius Thermal Damage ODE}

The thermal damage field $\Omega_d(\mathbf{r}, t)$ satisfies a
first-order integro-differential equation at each voxel:
\begin{equation}
    \frac{d\Omega_d}{dt}
    = A\,\exp\!\left(\frac{-E_a}{R\,T_K(\mathbf{r},t)}\right),
    \qquad \Omega_d(\mathbf{r}, 0) = 0
    \label{eq:arrhenius}
\end{equation}
where $T_K = T_{\text{celsius}} + 273.15$ is temperature in Kelvin,
$A = 3.1 \times 10^{98}\ \si{\per\second}$ is the frequency factor,
$E_a = 6.28 \times 10^5\ \si{\joule\per\mol}$ is the activation energy,
and $R = \SI{8.314}{\joule\per\mol\per\kelvin}$ is the gas constant.
The damage threshold $\Omega_d = 1$ corresponds to a 63\% probability
of irreversible cell necrosis (Henriques--Moritz criterion).

%----------------------------------------------------------------------
\section{Boundary Conditions}
\label{sec:bc}

Three boundary condition types are supported on $\partial\Omega$:
\begin{description}
    \item[Dirichlet] $T = T_{\mathrm{wall}}$ — models a large blood vessel
        or cooling pad holding the boundary at a fixed temperature
        (default: $T_{\mathrm{wall}} = T_b = \SI{37}{\degreeCelsius}$).
    \item[Neumann] $k\,\partial T/\partial n = 0$ — insulated / symmetry
        plane; built into the FD Laplacian via ghost-cell correction.
    \item[Robin] $k\,\partial T/\partial n + h_c(T - T_\infty) = 0$ —
        convective cooling at the tissue surface with heat transfer
        coefficient $h_c$ [\si{\watt\per\square\meter\per\kelvin}]
        and ambient temperature $T_\infty$.
\end{description}

%----------------------------------------------------------------------
\section{Cost Functional (Bolza Form)}

The optimal control problem minimizes:
\begin{equation}
    J[P(\cdot),\, t_f]
    = \underbrace{
        \gamma_1 \int_{\Omega_T}\!\!\max\bigl(0,\,1-\Omega_d(t_f)\bigr)^2\,d\mathbf{r}
        + \gamma_2\,t_f
      }_{\text{terminal cost } \Phi}
    + \int_0^{t_f}\!\!
      \underbrace{
        \left[
          \alpha_1\|P\|^2
          + \alpha_2 \int_{\Omega_H}\!\!\max(0,\,T - T_{\text{safe}})^2\,d\mathbf{r}
          + \alpha_3
        \right]
      }_{\text{running cost } L}\,dt
    \label{eq:cost}
\end{equation}
subject to~\eqref{eq:bioheat-discrete} and~\eqref{eq:arrhenius},
with $P(t) \in [P_{\min},\, P_{\max}]$ and $|\dot{P}| \leq \dot{P}_{\max}$.

Table~\ref{tab:cost-weights} lists the default weight values.

\begin{table}[htbp]
\centering
\caption{Default Bolza cost weights}
\label{tab:cost-weights}
\renewcommand{\arraystretch}{1.4}
\begin{tabular}{lll}
\toprule
Symbol & Role & Default \\
\midrule
$\alpha_1$ & Energy penalty ($\|P\|^2$) & $10^{-4}$ \\
$\alpha_2$ & Healthy-tissue overheating & $1.0$ \\
$\alpha_3$ & Time-rate cost (encourages speed) & $0.01$ \\
$\gamma_1$ & Incomplete ablation penalty & $10.0$ \\
$\gamma_2$ & Soft minimum-time penalty & $0.1$ \\
\bottomrule
\end{tabular}
\end{table}

%----------------------------------------------------------------------
\section{Adjoint Equations and Optimal Control Law}

Applying the \gls{pmp} to the discretized system, the costate vector
$\boldsymbol{\lambda} = [\boldsymbol{\lambda}_T,\, \boldsymbol{\lambda}_\Omega]^\top$
satisfies backward ODEs integrated from $t_f$ to $0$:
\begin{align}
    -\dot{\boldsymbol{\lambda}}_T
    &= A_d^\top\,\boldsymbol{\lambda}_T
       + 2\alpha_2\,\max(0,\, T - T_{\text{safe}})\cdot\mathbf{1}_{\Omega_H}
       \nonumber \\
    &\quad + \boldsymbol{\lambda}_\Omega \odot
       \left[A\,\frac{E_a}{R T_K^2}\,
             \exp\!\left(\frac{-E_a}{R T_K}\right)\right]
    \label{eq:adjoint-T} \\[4pt]
    -\dot{\boldsymbol{\lambda}}_\Omega &= \mathbf{0}
    \quad\Rightarrow\quad
    \boldsymbol{\lambda}_\Omega(t) = \boldsymbol{\lambda}_\Omega(t_f) = \text{const.}
    \label{eq:adjoint-Omega}
\end{align}
with terminal conditions:
\begin{align}
    \boldsymbol{\lambda}_T(t_f) &= \mathbf{0} \\
    \boldsymbol{\lambda}_\Omega(t_f)
    &= -2\gamma_1\,\max\bigl(0,\,1 - \Omega_d(t_f)\bigr)\cdot\mathbf{1}_{\Omega_T}
\end{align}

The Pontryagin optimal control law (projection onto the control interval) is:
\begin{equation}
    P^*(t)
    = \clip\!\left(\frac{\boldsymbol{\lambda}_T^\top\,\mathbf{b}_P}{2\alpha_1},\;
                   P_{\min},\; P_{\max}\right)
    \label{eq:pontryagin}
\end{equation}
